\chapter{REFERENCIAL TEÓRICO}
Este capítulo constitui a base teórica do seu TCC, o objetivo é apresentar, analisar e sintetizar as ideias, conceitos e teorias dos autores que fundamentam a sua investigação e sustentam a sua intervenção pedagógica.

Diferente de uma linha do tempo, o referencial deve ser estruturado por afinidade temática. Agrupe o texto em subseções que correspondam aos principais conceitos do seu problema de pesquisa.

Para elaborar o referencial teórico é importante:
\begin{uemaball}
    \item analisar a evolução cronológica do assunto nos documentos anteriormente publicados;
    \item ater-se às contribuições mais relevantes diretamente relacionadas ao assunto e utilizá-las na discussão dos resultados;
    \item metodologia ou materiais e métodos -- neste item deve-se considerar a descrição precisa dos métodos, técnicas e processos utilizados durante a pesquisa, proporcionando a repetição dos experimentos ou estudo por outros pesquisadores;
    \item resultados e discussão -- neste item devem ser descritos de forma objetiva os resultados das análises quantitativa e qualitativa dos dados e informações obtidas na pesquisa, confrontando ao problema, aos objetivos, às hipóteses e ao referencial teórico.
\end{uemaball}

Para garantir o rigor científico e evitar o plágio acadêmico, toda ideia que não seja de sua autoria deve ser obrigatoriamente referenciada. Exemplo de citação indireta \cite{towfic2022benchmarking}. Segundo \citeonline{chu2021sustained}, isto é um outro exemplo de citação no texto.

\begin{citacao}
Injeção de dependências (\textit{dependency injection} ou DI) é um padrão de
desenvolvimento de software usado para manter o baixo acoplamento entre classes do
sistema. As dependências de um objeto não são instanciadas programaticamente, mas
sim injetadas de alguma forma \cite{chambers2021macbook}.
\end{citacao}
